\subsection{How the Church Came to Keep Advent}\label{sec:history}

The weeks before Christmas were kept in different ways in different churches: in Gaul by a fast
from the feast of St Martin, at Rome by a fast in the tenth month and then by a set of Sunday
Masses, at Milan by six Sundays and in the old Spanish books by five. The books of 1962 inherit
the Roman form, with its four Sundays and its December Ember Days. The medieval commentators,
writing on the liturgy of their own day, explained the marks that set the season apart.

\subsubsection{In Spain and Gaul}

A council at Saragossa in Spain legislated for the days around Christmas. Hefele, the historian of
the councils, places the synod ``somewhere about the year 380,'' and notes that Mansi argued for
379; its eight canons are directed against the Priscillianists. The fourth reads, in Hefele's
English:

\begin{quote}
(4) From the 17th December to the Feast of the Epiphany every one must attend the church daily,
and may not go with bare feet.\par
{\smallskip\noindent\small ---Council of Saragossa, c.~4, in C.~J. Hefele, \textit{A History of
the Councils of the Church}, vol.~II (Edinburgh, 1876), pp.~292--293; Latin in Mansi III,
coll.~633--635.\par}
\end{quote}

The canon requires daily attendance at church through the days before and after Christmas. It
appoints no fast.

The fast before Christmas appears at Tours. Gregory of Tours, in the chapter of his
\textit{History of the Franks} that names the bishops of his see (X.31), says of Perpetuus, the
sixth of them, that he pulled down the basilica Bishop Bricius had built over St Martin and
built a larger one, into whose apse he moved the saint's body. Of the same bishop he
writes \latin{instituit ieiunia vigiliasque, qualiter per circulum anni observarentur} (gloss:
he instituted fasts and vigils, and how they were to be kept through the round of the year), a
rule which, Gregory adds, \latin{hodieque apud nos tenetur scriptum} (gloss: is kept in writing
among us to this day). He then copies the list. In Dalton's English it reads in part:

\begin{quote}
Fasts. After Pentecost, the fourth and sixth days of the week to the Nativity of Saint John. From
the first of September to the first of October, two fasts each week. From the first of October
to the burial of our lord Martin, two fasts each week. From the burial of our lord Martin to the
Nativity of our Lord, three fasts each week. From the Nativity of the holy Hilary to the middle
of February, two fasts each week.\par
{\smallskip\noindent\small ---Gregory of Tours, \textit{History of the Franks} X.31, tr. O.~M.
Dalton, vol.~II (Oxford, 1927), p.~472; Latin ed. W. Arndt, MGH SS rer. Merov. I (Hannover,
1884), pp.~444--445.\par}
\end{quote}

Dalton's note dates the burial of St Martin to 11 November. The Latin of the line that concerns
Advent is \latin{De depositione domni Martini usque natale Domini terna in septimana ieiunia}
(gloss: from the burial of the lord Martin to the Nativity of the Lord, three fasts in the
week). Not every manuscript carries it: the Paris editors of 1838 note that the article is not in
the codex Regius.

A council at Mâcon, summoned by King Guntram and attended, Hefele says, by twenty-one bishops
from several provinces, made a fast of the same pattern a canon:

\begin{quote}
9. From S. Martin's Day until Christmas, a fast must be kept on Monday, Wednesday, and Friday of
each week (sabbati, cf. Du Cange, s.v.). The Sacrifice must be offered after the manner of the
Lenten season. The canons, too, shall be read at this time, so that no one may plead ignorance
in defence of a fault.\par
{\smallskip\noindent\small ---First Council of Mâcon, c.~9, in Hefele, \textit{History of the
Councils}, vol.~IV (Edinburgh, 1895), pp.~403--404; Latin in F. Maassen, \textit{Concilia aevi
merovingici}, MGH Conc. I (Hannover, 1893), p.~157.\par}
\end{quote}

The date of this council is disputed. Hefele dates it 581; Guéranger, in his \textit{Liturgical
Year}, says it was ``held in 582''; Maassen's edition heads it 583. Maassen's note calls the canon's order the same
that Perpetuus of Tours had already set, and quotes Gregory's line about the fast from the burial
of St Martin. The canon gives the weeks from St Martin's Day to Christmas three fast days a week
and the Lenten order at the altar. Guéranger, reading Perpetuus's rule, drew the comparison
himself and was careful about what could be known:

\begin{quote}
It would be impossible to decide whether St. Perpetuus, by this regulation, established a new
custom, or merely enforced an already existing law. Let us, however, note this interval of forty,
or rather of forty-three days, so expressly mentioned, and consecrated to penance, as though it
were a second Lent, though less strict and severe than that which precedes Easter.\par
{\smallskip\noindent\small ---Prosper Guéranger, \textit{The Liturgical Year: Advent}, tr.
Laurence Shepherd (Dublin, 1870), vol.~I, pp.~24--25.\par}
\end{quote}

A council at Tours in 567 also speaks of fasting in December, but its rule is for monks.
Hefele's English of its canon (his c.~17; Maassen numbers it 18) sets out the monastic year and
ends: ``In September, October, and November, again, the fasts must be three times a week, as
before; but in December, until Christmas, daily. From Christmas to Epiphany there shall be daily
prandium, because every day is a festival.'' It binds the monasteries, not the faithful at large.

\subsubsection{The December fast at Rome}

Rome's preparation for Christmas grew from another root. Pope St Leo the Great preached a set of
sermons on the fast of the tenth month, December, which the Roman church kept as one of four
seasonal fasts:

\begin{quote}
Accordingly we keep the spring fast in Lent, the summer fast at Whitsuntide, the autumn fast in
the seventh month, and the winter fast in this which is the tenth month, knowing that there is
nothing unconnected with the Divine commands, and that all the elements serve the Word of God to
our instruction, so that from the very hinges on which the world turns, as if by four gospels we
learn unceasingly what to preach and what to do.\par
{\smallskip\noindent\small ---Leo the Great, \textit{Sermo} XIX.2, tr. C.~L. Feltoe, NPNF, 2nd
ser., vol.~12 (New York, 1895).\par}
\end{quote}

Leo appoints the days of the fast at the end of his sermons in a fixed form: ``On Wednesday and
Friday therefore let us fast: and on Saturday let us keep vigil with the most blessed Apostle
Peter\dots'' (\textit{Sermo} XII.4, tr. Feltoe). The days he names are those of the December Ember
Days, which the 1962 calendar keeps in the week after the Third Sunday of Advent, and the
Saturday's station is still St Peter's. Leo preaches them as the fast of the tenth month. The
name ``Advent'' does not occur in these sermons; where the English of the Breviary or of the
Nicene and Post-Nicene Fathers puts ``Advent'' into a heading or a lesson, the word is the
translator's. Section~\ref{sec:ember} returns to what Leo taught about the fast.

\subsubsection{Five weeks and four}

The Roman books show the season's Sunday Masses taking shape. The Gelasian sacramentary, in
H.~A. Wilson's edition, has five Masses for Advent. The Gregorian sacramentary known as the
Hadrianum, in Wilson's edition of its ninth-century manuscripts, has four Sundays of Advent, and
the fourth is marked \latin{vacat} (gloss: it is empty), without a Mass of its own. Rome, then,
had five and later four. Milan kept six: the Ambrosian book edited by Magistretti has six Sundays
of Advent, and its season begins after St Martin. Spain stands between. Férotin, editing the
Mozarabic sacramentary, notes that the printed missal has six Sundays of Advent but that the old
manuscripts never have more than five.

The medieval commentators knew both numbers, and each accounted for the four weeks in his own
way. Amalarius of Metz found five weeks before Christmas marked in old lectionaries and four in
the Missal and the Antiphoner, and he traced the four to the Missal called Gregorian, of which he
says in another work \latin{cujus auctorem credimus esse beatum papam Gregorium} (gloss: whose
author we believe to be the blessed Pope Gregory). His word is \latin{credimus}, ``we believe'':
the ascription of the four-week arrangement to Gregory the Great comes to us as a report of
belief. Sicard of Cremona still saw both numbers in the books of his day, some marking five weeks
and others four. Durandus, too, knew old graduals and lectionaries that carry a \latin{Quinta
dominica ante natiuitatem} (gloss: fifth Sunday before the Nativity), but he gives the four weeks
an older author. He states, in an attribution this book does not take as history, that St Peter
first instituted three full weeks before Christmas and a fourth that is never completed, because
the Vigil of the Nativity is not of Advent.

Gregory the Great's own homilies point the same way. His headings, in the Latin text of his
homilies on the Gospels, place the homily on Luke 21 on the second Sunday of Advent of his day,
the homily on Matthew 11 on the third, and the homily on John 1 on the fourth; the homily on
Luke 3 he preached on the Ember Saturday. The 1962 Office reads the openings of the same homilies
on the First, Second, Third and Fourth Sundays.

\subsubsection{What the medieval commentators saw}

The commentators knew the season first by what it leaves out, and the list grew as they wrote.
Amalarius of Metz saw the \latin{Gloria in excelsis} omitted in Advent and, in some places, the
dalmatic laid aside. Rupert of Deutz adds that the dismissal \latin{Ite, missa est} is omitted.
John Beleth puts Advent beside Septuagesima, both without the \latin{Te Deum laudamus}, the
\latin{Gloria in excelsis} and the \latin{Ite, missa est}, and then marks where they part: in
Advent the \latin{Alleluia} is kept, because this time is partly a time of joy, while
Septuagesima is wholly a time of sorrow.

They also give reasons. Honorius Augustodunensis reads the four weeks as the time of the Old
Testament: the \latin{Gloria} and the dalmatic are set aside \latin{ad insinuandum tempus Veteris
Testamenti} (gloss: to signify the time of the Old Testament), which ran its course under four
orders that foreshadowed Christ's birth, the patriarchs, the judges, the kings and the prophets.
Durandus says that the \latin{Gloria} falls silent so that it may be sung more eagerly on the
night of the Nativity. He takes it for granted that Advent is a fast and that weddings are not
celebrated in it, and he asks why, in a fast, the \latin{Alleluia} is not silenced too. His answer
is Beleth's: this fast is partly of exultation and partly of mourning, and so it keeps one song
of joy and gives up another.

\subsubsection{The closed time}

Durandus's remark about weddings reflects an old discipline. The Council of Trent kept it: ``The
holy Synod enjoins that the ancient prohibitions of solemn nuptials be carefully observed by all,
from the Advent of our Lord Jesus Christ until the day of the Epiphany, and from Ash-Wednesday
until the octave of Easter inclusively'' (Session XXIV, \textit{De reformatione matrimonii},
c.~10, tr. J. Waterworth, p.~204). The Baltimore \textit{Manual of Prayers} of 1889, printed for
Catholics in the United States, names the time from Advent Sunday until the Epiphany as one in
which marriages are not solemnized.

The Code of Canon Law of 1917, the law of the Latin Church from 19 May 1918 to 26 November 1983,
narrowed the rule. Canon 1108 \S\,2 reads: \latin{Sollemnis tantum nuptiarum benedictio vetatur
a prima dominica Adventus usque ad diem Nativitatis Domini inclusive} (gloss: only the solemn
nuptial blessing is forbidden from the First Sunday of Advent to the day of the Nativity
inclusive).

\subsubsection{The season in the books of 1962}

The Code of Rubrics of 1960, in force from 1 January 1961 and followed by the books of 1962,
fixes the season and its days. Advent runs from First Vespers of the First Sunday to None of the
Vigil of the Nativity inclusive (n.~71). The four Sundays are of the first class (n.~11). The
ferias from 17 to 23 December and the Ember ferias are of the second class; the other weekdays of
Advent are of the third class and give way to feasts of the third class (nn.~24--25). Rose
vestments may be used on the Third Sunday of Advent, in its Office and Mass only (n.~131), and
the folded chasubles and the broad stole are no longer used (n.~137). Every Mass ends with
\latin{Ite, missa est}, and Advent is not among the exceptions the Code names (n.~507); the older
omission that Rupert and Beleth describe has gone. The instruction of the Sacred Congregation of
Rites on sacred music of 1958 keeps the instruments silent through the season, with the
exceptions Section~\ref{sec:third-sunday} sets out (nn.~81\,a, 83).

The Roman books do not describe the preparation for the Nativity in the Christian East.
Guéranger, writing in the nineteenth century, reported that the Greek Church of his
day kept a fast before Christmas of forty days from 14 November, the feast of the Apostle St
Philip, during which the people abstained from flesh-meat, butter, milk and eggs but were allowed
fish, oil and wine; fasting in the strict sense, he wrote, was binding on only seven of the forty
days, and the whole period went by the name of ``St. Philip's Lent.''
